\input{estilo}
\begin{document}
\begin{figura}\begin{tikzpicture}[arch/.style={archivo, minimum width=10mm, minimum height=13mm}]
  % el nodo funcionando como mapper
  \draw[marco] (0,0) rectangle (6.4,-6.75);
  \node[anchor=north west] at (0.15,-0.15) {mapper 1};
  \node[arch] (in) at (1.4,-1.6) {};
  \node[etiq, below=2pt] at (in.south) {input};
  \node[caja, minimum width=16mm, minimum height=9mm] (map) at (3.9,-1.6) {map};
  \draw[->] (in) -- (map);
  % la tira de pares pasa por el hash y el módulo
  \node[caja, minimum width=16mm, font=\scriptsize] (hash) at (3.9,-3.1) {hash mod \var{R}};
  \draw[->] (map) -- (hash);
  % los dos archivos en el disco local, uno por reducer
  % los archivos con su nombre, centrados en el recuadro con el mismo margen
  % en los cuatro lados: ocupan de y = -3.95 a y = -5.65
  \draw[marco, draw=suave] (1.6,-3.7) rectangle (6.2,-5.8);
  \node[etiq, below=1pt] at (3.9,-5.8) {disco local};
  \node[arch] (f1) at (2.6,-4.6) {};
  \node[arch] (f2) at (5.2,-4.6) {};
  \node[cod, below=2pt] at (f1.south) {M1\_R1};
  \node[cod, below=2pt] at (f2.south) {M1\_R2};
  \draw[->] (hash) -- (f1);
  \draw[->] (hash) -- (f2);
  % y cada uno, eventualmente, a su reducer
  \foreach \i/\x in {1/2.6, 2/5.2} {
    \node[caja, minimum width=11mm, minimum height=10mm] (r\i) at (\x,-7.9) {R\textsubscript{\i}};
    \draw[->, dashed] (\x,-5.65) -- (r\i.north);
  }
\end{tikzpicture}\end{figura}
\end{document}
